\chapter{Experiments and Results}
\label{ch:experiments}

TODO: in chapter 5 we should make deltaxGE the king of the evaluation.

TODO throughout:
- reference eval equivalent
- if in ch5 some formulas are missing we have to add them. in particular all plot lines and all values in the tables

This chapter applies the evaluation framework of \Cref{ch:eval} to characterize
the proposed guidance method. We first inspect representative guided samples
qualitatively, then compare alternative flow-time guidance profiles, and finally
study the effect of varying the target specification.


\section{Experimental setup}
\label{sec:experimental-setup}

\subsection{Experimental settings}
\label{sec:experimental-settings}

All experiments consider a single generated weather state, as specified in
\Cref{ch:eval}. We denote an experimental setting by

\begin{equation}
    e = (\tau,\rho_H,\pi),
    \label{eq:experimental-setting}
\end{equation}

where $\tau$ is the forecast initialization, $\rho_H$ the prescribed relative
target intensity, and $\pi$ the spatial target mask. The target variable and
level are fixed throughout the experiments unless stated otherwise.

For the guidance-profile comparison in T1, we use a fixed target region over
Europe and the four forecast initializations

\[
\begin{aligned}
    &\text{9 August 2024, 12:00 UTC}, \\
    &\text{12 August 2024, 12:00 UTC}, \\
    &\text{4 July 2026, 12:00 UTC}, \\
    &\text{9 July 2026, 12:00 UTC}.
\end{aligned}
\]

Each initialization is evaluated at

TODO: it would be nice to have the implied temperature for each experiment.
Or if not we can in the qualitative analysis show a table statistic with the min,max,mean,std.

\begin{equation}
    \rho_H \in \{0.002,\,0.004\},
    \label{eq:t1-target-intensities}
\end{equation}

corresponding to relative target increases of $0.2\%$ and $0.4\%$.
Since the mask is fixed throughout T1, we abbreviate these settings as
$e=(\tau,\rho_H)$. The four initializations and two intensities yield eight
experimental settings.

For each setting, we generate $M=5$ stochastic ensemble members indexed by

\begin{equation}
    m \in \{0,\ldots,M-1\}.
\end{equation}

Each candidate guidance profile is therefore evaluated on $8\times5=40$
paired realizations.

\begin{figure}[H]
    \centering
    \includegraphics[width=0.65\textwidth]{figures/mask_flat.png}
    \caption{
        Spatial target mask used throughout the T1 guidance-profile comparison.
    }
    \label{fig:t1-target-mask}
\end{figure}


\subsection{Paired stochastic realizations}
\label{sec:experimental-seeding}

Comparisons between guidance configurations use matched stochastic
realizations. The random seed is determined from the weather-step and
ensemble-member indices as

\begin{equation}
    \mathrm{seed}(n,m)=1000n+m.
    \label{eq:experimental-seed}
\end{equation}

The seed is kept fixed across configurations that are compared directly.
Corresponding guided and unguided states therefore originate from the same
initial residual-noise realization. This controls sampling noise in paired
comparisons while retaining stochastic variability across ensemble members.


\subsection{Guidance-profile search}
\label{sec:guidance-profile-search}

We determine the flow-time guidance profile in two stages. First, we test how
the response depends on the flow step at which guidance is concentrated using

\begin{equation}
    \Gamma_{\mathrm{spike}}
    =
    \left\{
        \mathrm{spike@0},
        \mathrm{spike@1},
        \mathrm{spike@2},
        \mathrm{spike@5},
        \mathrm{spike@11},
        \mathrm{spike@17},
        \mathrm{spike@23}
    \right\}.
    \label{eq:spike-profile-set}
\end{equation}

Here, $\mathrm{spike@}k$ assigns unit relative guidance strength to flow step
$k$ and zero to all remaining steps.

In a second experiment, the strongest localized candidate is compared with
profiles that distribute the guidance signal over neighbouring flow steps. A
profile $\mathrm{spread@}a\text{-}b$ decreases linearly from one at step $a$
to zero at step $b$. For example,

\[
    \mathrm{spread@0\text{-}2}
    =
    (1,\;0.5,\;0)
\]

over flow steps $0$, $1$, and $2$. The specific spread profiles considered are
introduced together with the corresponding experiment in
\Cref{sec:experiments-temporal-spreading}.


\subsection{Implementation details}
\label{sec:experimental-implementation}

All experiments use the same numerical sampling and guidance implementation.
The residual flow is discretized into 25 flow steps.

% TODO: Add secant-search tolerance and maximum number of iterations.
% TODO: Add numerical precision used during inference.
% TODO: Add compute hardware.

TODO: insert rest once chapter done


\section{T0: Qualitative evaluation}
\label{sec:experiments-qualitative}

We begin with a qualitative inspection of representative guided samples.
Unless stated otherwise, we use

(TODO: reference only date in format 2024-08-09, time is fixed to 12:00 and we should say that in the setup)

\begin{equation}
    e_{\mathrm{ex}}
    =
    (\text{9 July 2026, 12:00 UTC},\,0.002),
    \qquad m=0.
    \label{eq:representative-experiment}
\end{equation}


\subsection{Descriptive statistics}

We start with a dry factual description of what happened.
Put ung on left and gui on right for each metric in a split row.
Put here table with min,max,mean,std for each experiment and rho (rhos tables with exps rows). 


\subsection{Guided and matched unguided states}

In the following we show the t2m field for the ung | gui and gui states with the corresponding guidance effect (the difference between the two objects).
Figure~\ref{fig:t0-state-comparison} illustrates the three objects used
throughout the evaluation: the matched unguided state
$x^{\mathrm{ung}\mid\mathrm{gui}}$, the guided state $x^{\mathrm{gui}}$, and
their difference $\Delta x^{\mathrm{GE}}$.

This shows us where the guidance effect has taken place, and how the guided state looks like when it's applied over the ung|gui twin.

The top of the charts also report descriptive statistics to orient oneself.

%%%%%%%%%%%%
TODO: fill the page in a smarter way

\begin{figure}[H]
    \centering

    % Top row: ung|gui and GE
    \includegraphics[
        width=\textwidth,
        height=0.2\textheight,
        keepaspectratio
    ]{figures/ung_gui.png}
    \includegraphics[
        width=\textwidth,
        height=0.2\textheight,
        keepaspectratio
    ]{figures/gui.png}
    % Bottom row: gui
    \includegraphics[
        width=\textwidth,
        height=0.2\textheight,
        keepaspectratio
    ]{figures/GE.png}

    \caption{
        Matched unguided state $x^{\mathrm{ung}\mid\mathrm{gui}}$ (top left),
        guidance effect $\Delta x^{\mathrm{GE}}$ (top right), and guided state
        $x^{\mathrm{gui}}$ (bottom) for the representative experiment
        $e_{\mathrm{ex}}$ using $\mathrm{spike@1}$.
    }
    \label{fig:t0-state-comparison}
\end{figure}
%%%%%%%%%%%%

\subsection{Masked mean}

\begin{figure}[H]
    \centering
    \includegraphics[width=\textwidth]{figures/T0.0.png}

    \vspace{0.5em}

    \includegraphics[width=\textwidth]{figures/T0.1.png}

    \caption{
        ...
    }
    \label{fig:...}
\end{figure}

TODO: split figures and put after corresponding text chunks

First one shows the average value, to be able to understand the magnitude of the guidance effect.
(not so much visually from the plot, enables us to anchor to reality)

Second one shows where the effect has taken place in the levels (we can compare levels) in terms of the masked value wrt to the one of the ung|gui.


Here we should notice how guiding at the last step clearly doesn't produce any change in the other variables. Which can also be observed under the second "Evolution through the residual flow" plot and 
under the "Spatial structure of the guidance effect"

\subsection{Evolution through the residual flow}

Figure~\ref{fig:t00} shows the evolution of
$\bar z^{\mathrm{gui}}_{t,m,v}$ over flow time for the selected variables.
The curves show the ensemble mean, while the shaded region spans the minimum
and maximum across the five ensemble members. The shading enables to see the variance of the noise drawings (different starting points) and the less varying end points of the trajectories.

\begin{figure}[H]
    \centering
    \includegraphics[width=\textwidth]{figures/T00.png}
    \caption{
        Evolution of the masked guided residual
        $\bar z^{\mathrm{gui}}_{t,m,v}$ for selected atmospheric variables.
        Curves show the ensemble mean, while the shaded region spans the
        minimum and maximum across ensemble members. Upper-air variables are
        shown at their highest available pressure level for visualization.
    }
    \label{fig:t00}
\end{figure}


We next inspect how the intervention acts along the residual flow, by visualizing the same object but splitting guided and unguided steps along the trajectory.
Figure~\ref{fig:t0-early-late} compares an early intervention,
$\mathrm{spike@1}$, with the later $\mathrm{spike@17}$ profile under the same
experimental setting and stochastic realization.
(TODO: The plot also let you see visually that guidance is concentrated in one single step)

(TODO: the plot has 1-e5 over q ... remove)
\begin{figure}[H]
    \centering
    \includegraphics[width=\textwidth]{figures/T01a.png}

    \vspace{0.5em}

    \includegraphics[width=\textwidth]{figures/T01b.png}

    \caption{
        Step-wise decomposition of the guided residual-flow trajectory for
        $\mathrm{spike@1}$ (top) and $\mathrm{spike@17}$ (bottom), using the
        same experimental setting and ensemble realization. The panels show
        the current state, the corresponding
        unguided-under-guided-context landing, and the contributions of the
        flow and guidance moves at each flow step.
    }
    \label{fig:t0-early-late}
\end{figure}

To compare the residual trajectories across atmospheric variables, we apply
the spatial target mask to the guided residual state and define

TODO: simply rewrite with the mask operator and distinguish gui and ung step.


Across all experiments, all candidate spike profiles reach their
prescribed targets. The profiles therefore differ in the atmospheric response
through which the target is realized, which we compare quantitatively next.


\section{T1: Guidance-profile evaluation}
\label{sec:experiments-guidance-profiles}

We now compare the candidate guidance profiles using the quantitative criteria
defined in \Cref{sec:guidance-profile-evaluation}.


\subsection{T1.1: Temporal localization of guidance}
\label{sec:experiments-temporal-localization}

TODO: this is now the sample charts for exposing the quantity.

\begin{figure}[H]
    \centering
    \includegraphics[width=\textwidth]{figures/T02a.png}

    \vspace{0.5em}

    \includegraphics[width=\textwidth]{figures/T02b.png}

    \caption{
        Spatial guidance effect $\Delta x^{\mathrm{GE}}$ for 2\,m temperature
        (top) and geopotential at 1000\,hPa (bottom) in the representative
        experiment. The response in geopotential illustrates how an
        intervention targeted at 2\,m temperature propagates to another
        atmospheric variable.
    }
    \label{fig:t0-spatial-effect}
\end{figure}

We first compare the seven profiles in
$\Gamma_{\mathrm{spike}}$, which differ only in the flow step at which guidance
is applied.

Figure~\ref{fig:t1-propagation} shows the resulting multivariate propagation.
The response peaks at $\mathrm{spike@1}$ and decreases strongly as guidance is
shifted toward later flow steps.
(TODO: this is correct, but there is also spike@0 which before spike@1 that is less. same argument goes for the next two tests)

(TODO: this is the same chart as the second of the masked average subsection just in terms of absolute values so we can see how much levels where pushed respectively, instead of the direction)
\begin{figure}[H]
    \centering
    \includegraphics[width=\textwidth]{figures/T1.png}
    \caption{
        Multivariate propagation of the guidance effect for the candidate
        spike profiles. Earlier interventions, and in particular
        $\mathrm{spike@1}$, produce a stronger response across atmospheric
        variables and pressure levels.
    }
    \label{fig:t1-propagation}
\end{figure}

TODO: reference the criteria from latest chapter (this is something we have to do throughout)
The spatial-deviation criterion exhibits the same preference for early
guidance. Early interventions depart more strongly from the prescribed mask,
whereas later interventions increasingly reproduce its spatial structure.

TODO: should also add the z1000 one.
\begin{figure}[H]
    \centering
    \includegraphics[width=\textwidth]{figures/T2.1.png}
    \caption{
        Spatial deviation of the normalized guidance footprint from the
        prescribed target mask, measured by total-variation distance. Larger
        values indicate a stronger departure from the imposed mask geometry.
    }
    \label{fig:t1-spatial-deviation}
\end{figure}

Ensemble variability follows the same pattern. It is largest for
$\mathrm{spike@1}$ and decreases toward later flow steps, indicating that
early guidance retains a stronger dependence on the stochastic realization.

TODO: should also add the z1000 one.
\begin{figure}[H]
    \centering
    \includegraphics[width=\textwidth]{figures/T3.1.png}
    \caption{
        Ensemble variability of the guidance footprint for the candidate spike
        profiles. Larger values indicate a stronger dependence of the induced
        perturbation on the stochastic ensemble realization.
    }
    \label{fig:t1-ensemble-variability}
\end{figure}

The three criteria are summarized in
\Cref{tab:t1-spike-summary}. 
(TODO: this has to be clarified)
The propagation and spatial-deviation statistics
are pooled across the 40 setting-member combinations, while ensemble
variability is first computed across the five ensemble members within each
experimental setting and then summarized across settings and target
intensities.

\begin{table}[H]
    \centering
    \caption{
        Summary of the temporal-localization experiment across the four
        initializations and two target intensities. Larger values correspond
        to the preferred behaviour for all three evaluation criteria.
    }
    \label{tab:t1-spike-summary}
    \begin{tabular}{lccc}
        \toprule
        Profile
        & Multivariate propagation
        & Spatial deviation
        & Ensemble variability \\
        \midrule
        $\mathrm{spike@0}$
        & $0.680\pm0.043$
        & $0.337\pm0.046$
        & $0.094\pm0.030$ \\

        $\mathrm{spike@1}$
        & $\mathbf{1.000\pm0.098}$
        & $\mathbf{0.357\pm0.053}$
        & $\mathbf{0.108\pm0.035}$ \\

        $\mathrm{spike@2}$
        & $0.759\pm0.053$
        & $0.341\pm0.049$
        & $0.098\pm0.033$ \\

        $\mathrm{spike@5}$
        & $0.454\pm0.018$
        & $0.290\pm0.035$
        & $0.081\pm0.027$ \\

        $\mathrm{spike@11}$
        & $0.130\pm0.006$
        & $0.191\pm0.018$
        & $0.056\pm0.019$ \\

        $\mathrm{spike@17}$
        & $0.044\pm0.001$
        & $0.093\pm0.007$
        & $0.027\pm0.008$ \\

        $\mathrm{spike@23}$
        & $0.027\pm0.001$
        & $0.008\pm0.001$
        & $0.003\pm0.001$ \\
        \bottomrule
    \end{tabular}
\end{table}

Taken together, the three criteria give a consistent picture. Early guidance
induces a stronger multivariate response, departs more substantially from the
prescribed spatial mask, and retains greater variability across ensemble
members. All three effects are strongest for $\mathrm{spike@1}$ and diminish
as guidance is moved toward later flow steps. We therefore retain
$\mathrm{spike@1}$ as the localized reference profile.

\subsection{T1.2: Temporal spreading of guidance}
\label{sec:experiments-temporal-spreading}

Having identified $\mathrm{spike@1}$ as the strongest localized profile, we
next test whether distributing guidance over neighbouring flow steps improves
the resulting response. We compare it with

\begin{equation}
    \Gamma_{\mathrm{spread}}
    =
    \left\{
        \mathrm{spread@0\text{-}2},
        \mathrm{spread@1\text{-}3}
    \right\},
    \label{eq:spread-profile-set}
\end{equation}

using the same experimental settings and paired stochastic realizations as in
T1.1. The pooled results are summarized in
\Cref{tab:t1-spread-summary}.

\begin{table}[H]
    \centering
    \caption{
        Comparison of $\mathrm{spike@1}$ with guidance distributed over
        neighbouring flow steps, summarized across the four initializations
        and two target intensities. Larger values correspond to the preferred
        behaviour for all three criteria.
    }
    \label{tab:t1-spread-summary}
    \begin{tabular}{lccc}
        \toprule
        Profile
        & Multivariate propagation
        & Spatial deviation
        & Ensemble variability \\
        \midrule
        $\mathrm{spike@1}$
        & $\mathbf{0.999\pm0.098}$
        & $0.357\pm0.053$
        & $\mathbf{0.1080\pm0.035}$ \\

        $\mathrm{spread@0\text{-}2}$
        & $0.746\pm0.045$
        & $0.345\pm0.047$
        & $0.0980\pm0.032$ \\

        $\mathrm{spread@1\text{-}3}$
        & $0.949\pm0.073$
        & $\mathbf{0.359\pm0.050}$
        & $0.1077\pm0.036$ \\
        \bottomrule
    \end{tabular}
\end{table}

Spreading guidance over steps $0$--$2$ reduces all three scores.
For $\mathrm{spread@1\text{-}3}$, the differences are smaller:
multivariate propagation and ensemble variability decrease, while spatial
deviation increases marginally. Thus, distributing guidance over neighbouring
flow steps does not provide a consistent improvement over the localized
$\mathrm{spike@1}$ intervention.

\subsection{T1.3: Profile selection}
\label{sec:experiments-profile-selection}

Taken together, the temporal-localization and spreading experiments support
retaining $\mathrm{spike@1}$. It achieves the strongest response among the
localized profiles, while distributing guidance over neighbouring flow steps
does not provide a consistent improvement. In particular,
$\mathrm{spread@1\text{-}3}$ performs essentially on par in spatial deviation
and ensemble variability but exhibits weaker multivariate propagation.

Given these comparable results, we favour the simpler localized intervention
and select

\begin{equation}
    \gamma^\star = \mathrm{spike@1}
    \label{eq:selected-guidance-profile}
\end{equation}

for the remaining experiments.

\subsection{T1.4: Prediction accuracy over flow time}
\label{sec:experiments-prediction-noise}

The preference for $\mathrm{spike@1}$ does not coincide with increased
accuracy of the intermediate residual prediction. To examine this, we compute

\begin{equation}
    \operatorname{RMSE}\!\left(\hat z_t,z_T\right)
\end{equation}

between the intermediate prediction $\hat z_t$ and the final residual $z_T$.

\begin{figure}[H]
    \centering
    \includegraphics[width=\textwidth]{figures/T4.png}
    \caption{
        RMSE between the intermediate residual prediction $\hat z_t$ and the
        final residual $z_T$ over flow time. Curves summarize the four forecast
        initializations, with variation across ensemble members.
    }
    \label{fig:t1-prediction-noise}
\end{figure}

For all four initializations, the prediction error reaches its maximum at
$t=1$, precisely the flow step selected for guidance. Thereafter, the error
decreases approximately monotonically toward the final residual. Thus, the
effectiveness of $\mathrm{spike@1}$ cannot be explained by unusually accurate
intermediate predictions. Instead, the relationship between prediction
accuracy and guidance effectiveness is non-monotonic and remains unresolved by
the present experiments.





% \section{T2: Scenario-specification analysis}
% \label{sec:experiments-scenario-specification}

% Having selected
% $\gamma^\star=\mathrm{spike@1}$, we now examine how the guidance response
% changes with the prescribed scenario. We vary the target intensity and spatial
% extent while keeping the remaining experimental configuration fixed.


% \subsection{T2.1: Target intensity}
% \label{sec:experiments-target-intensity}

% We first vary the prescribed relative intensity $\rho_H$ while keeping the
% target mask and guidance profile fixed. This tests how increasingly extreme
% targets affect the magnitude and spatial structure of the required
% intervention.

% % TODO: State the set of target intensities used in this experiment.

% % \begin{figure}[H]
% %     \centering
% %     % TODO: Insert target-intensity figure.
% %     \includegraphics[width=\textwidth]{figures/TODO_target_intensity.png}
% %     \caption{
% %         Guidance response for increasing prescribed target intensities using
% %         the selected profile $\gamma^\star=\mathrm{spike@1}$. The comparison
% %         shows how the generated state and guidance effect evolve as the target
% %         is moved progressively further from the matched unguided reference.
% %     }
% %     \label{fig:t2-target-intensity}
% % \end{figure}

% % TODO: Add the result once the target-intensity experiment is summarized.

% \subsection{T2.2: Target region}
% \label{sec:experiments-target-region}

% We next vary the spatial extent of the target mask while keeping the target
% variable, target intensity, and guidance profile fixed. The masks share the
% same center but differ in spatial extent.

% % TODO: State the set of mask sizes used in this experiment.

% % \begin{figure}[H]
% %     \centering
% %     % TODO: Insert target-region figure.
% %     \includegraphics[width=\textwidth]{figures/TODO_target_region.png}
% %     \caption{
% %         Guidance response for target regions of increasing spatial extent using
% %         $\gamma^\star=\mathrm{spike@1}$. The comparison shows how the generated
% %         state and guidance effect change as the prescribed intervention covers
% %         a larger region.
% %     }
% %     \label{fig:t2-target-region}
% % \end{figure}

% % TODO: Add the result once the target-region experiment is summarized.

% \section{Chapter summary}
% \label{sec:experiments-summary}

% The experiments identify $\mathrm{spike@1}$ as the preferred guidance profile.
% Guidance applied at this flow step produces the strongest multivariate
% propagation, the largest departure from a direct imprint of the prescribed
% mask, and the strongest dependence on the stochastic ensemble realization.
% Spreading the intervention over neighbouring flow steps does not improve these
% properties.

% Interestingly, the selected flow step also coincides with the largest
% intermediate residual-prediction error, indicating that guidance effectiveness
% is not simply determined by prediction accuracy. The T2 experiments further
% examine how the selected configuration responds to changes in target intensity
% and spatial extent.